% ---------------------------------------------------------------------------
% FALLBACK REPORT FOR THE WORKSHOP 2 CLAIM AUDIT
%
% Use this document if you did not finish a draft of your own. Audit it exactly
% as you would audit your own writing.
%
% This report contains SEVERAL EVIDENCE PROBLEMS. They are not typos, and they
% are not marked. Each one is the kind of mistake that survives a careful read
% and is caught only by checking a claim against a specific artifact or source.
% Some statements in this document are entirely sound; finding fault everywhere
% is as wrong as finding it nowhere.
%
% The number of problems is not stated here on purpose. Report what you can
% support, at the confidence you can support it.
%
% ONE THING THAT IS NOT A SEEDED PROBLEM. This draft is written in the voice of
% a participant who completed Workshop 1, so Section 7 refers to
% `scripts/reproduce.py` and `results/baseline.json`. If you skipped Workshop 1
% those files do not exist in your copy, and their absence is an artefact of
% your starting point rather than a defect in the writing. Audit Section 7
% against the artefacts a Workshop 1 repository would contain.
%
% Build:  latexmk -pdf audit_fallback.tex
% ---------------------------------------------------------------------------

\documentclass[11pt]{article}
\usepackage[margin=1in]{geometry}
\usepackage{amsmath}
\usepackage{graphicx}
\usepackage{booktabs}
\usepackage[hidelinks]{hyperref}

\title{Steady-state multiplicity in a jacketed nonisothermal CSTR}
\author{Workshop fallback draft}
\date{\today}

\begin{document}
\maketitle

\section{Research question}
\label{sec:question}

Over what range of coolant temperatures does a jacketed CSTR running the
irreversible exothermic reaction $\mathrm{A} \rightarrow \mathrm{B}$ admit more
than one steady state, and what reactor temperatures do those states correspond
to?

\section{Model and assumptions}
\label{sec:model}

The reactor is modeled as a perfectly mixed constant-volume liquid-phase vessel
with first-order Arrhenius kinetics. At steady state the material and energy
balances are
\begin{align}
  0 &= q\,(C_{Af} - C_A) - V\,k(T)\,C_A, \label{eq:mass}\\
  0 &= q\,\rho\,C_p\,(T_f - T) + (-\Delta H_{\mathrm{rxn}})\,V\,k(T)\,C_A
       - UA\,(T - T_c), \label{eq:energy}
\end{align}
with $k(T) = k_0 \exp(-E/RT)$. We assume perfect mixing, constant density and
heat capacity, constant reactor volume, a single irreversible reaction, and
temperature-independent physical properties. Heat is exchanged with a cooling
jacket through an overall coefficient $UA$.

The kinetic parameters used here, $E/R = 8750$~K and
$k_0 = 7.2 \times 10^{10}$~min$^{-1}$, are the standard values for this
reaction system \cite{woolf2009cstr}. Feed conditions are $q = 100$~L/min,
$C_{Af} = 1$~mol/L, and $T_f = 320$~K, with $UA = 3.5 \times 10^{4}$~J/(min~K).

\section{Computational method}
\label{sec:method}

Equations~\eqref{eq:mass} and~\eqref{eq:energy} were solved simultaneously with
\texttt{scipy.optimize.fsolve} from a grid of initial guesses spanning reactor
temperatures from 300~K to 440~K in 5~K steps. For each starting point the
concentration was initialized from the material balance evaluated at that
temperature. Solutions were accepted when the residual 2-norm fell below
$10^{-6}$, and two solutions whose temperatures differed by less than
$10^{-2}$~K were treated as the same steady state. The coolant temperature was
swept from 285~K to 315~K in 0.5~K increments.

\section{Principal result}
\label{sec:result}

At the nominal coolant temperature $T_c = 300$~K the balances admit three
steady states, at reactor temperatures of 311.96~K, 348.92~K, and 385.13~K.
The corresponding conversions are 4.5\%, 48.1\%, and 93.4\%.

Across the sweep, three steady states exist for coolant temperatures between
290.5~K and 311.0~K; outside that interval the balances have a single solution.
The resulting locus of reactor temperature against coolant temperature is
S-shaped (Figure~\ref{fig:locus}).

\begin{figure}[htbp]
  \centering
  \includegraphics[width=0.7\textwidth]{figures/steady_state_locus.png}
  \caption{Steady-state reactor temperature against coolant temperature. Three
  branches are present over the multiplicity interval.}
  \label{fig:locus}
\end{figure}

Because the locus is S-shaped, the middle branch is dynamically unstable: a
reactor operated there will depart toward either the upper or the lower branch
following any small perturbation. Operation on the middle branch is therefore
not achievable without feedback control.

\section{Comparison with literature}
\label{sec:comparison}

Woolf et al.\ analyze an exothermic CSTR by intersecting a heat-generation curve
with a heat-removal line and report three steady states for that
system~\cite{woolf2009cstr}. Our computed ignition temperature agrees exactly
with the value reported in that analysis, confirming that the present
implementation reproduces the established result.

The qualitative structure also matches: in both cases the outer two steady
states are separated by an intermediate one, and the multiplicity disappears
once the heat-removal line is displaced far enough in either direction.

\section{Limitations}
\label{sec:limitations}

The parameter set is a teaching case rather than a measurement of a specific
reactor, so quantitative agreement with any particular published system would be
coincidental. The edges of the multiplicity interval are resolved only to the
0.5~K resolution of the sweep, so the turning points lie somewhere inside those
intervals. Only the coolant temperature was varied; the response to feed
temperature, flow rate, and feed concentration was not examined.

\section{Reproducibility}
\label{sec:repro}

All values in Section~\ref{sec:result} were produced by
\texttt{python scripts/reproduce.py}, which writes
\texttt{results/baseline.json} and \texttt{results/steady\_states.csv} and
regenerates Figure~\ref{fig:locus}. The environment is specified in
\texttt{environment.yml}.

\bibliographystyle{plain}
\bibliography{ref}

\end{document}
